% !TEX root = ../main.tex
\section{讨论与展望}

\subsection{方法有效性与机制分析}

本文从一个具体的传播模型出发，但希望回答的是更一般的问题：当复杂系统只留下多变量时间序列时，如何把演化规律与交互结构重新分离出来。
SIR 的价值在于它足够简单，
使真实感染曲线、传播参数和网络边都可以被精确检查；
同时它又包含随机性、非线性交互和吸收态，
能够暴露普通分类指标不易发现的问题。

实验首先说明，准确的一步预测并不等同于可信的动力学学习。
均值聚合和单向边表在常见状态上仍能得到一定预测能力，却难以稳定恢复感染概率对邻居数量的响应。
把自身演化与邻居交互分开表示后，
模型在不同网络结构和多步推演中的校准更稳定。
这里真正可推广的不是“感染者邻居计数”本身，
而是根据系统机理区分内在动力与外部耦合，
并选择能够保留交互强度的聚合方式。

实验还说明，动力学预测与拓扑解释必须相互验证。
未知网络时，稠密错误图仍可能通过动力学参数补偿获得看似可接受的短期输出。
硬图训练减少了一类连续缩放补偿，
分层样本让更多结构事件进入训练，
多初始化则帮助避开部分局部解；
然而九次运行中的两次严重坍缩表明，
结构辨识并没有因这些设计而被彻底解决。

\subsection{方法的可扩展性}

双通道框架依赖的核心结构是\cref{eq:local-interaction}，而不是 SIR 的具体概率公式。
只要组分状态能够被观测，
节点更新可以近似拆分为自身演化与邻居作用，
方法就具有迁移潜力。
例如，
SIS、SEIR 等传播模型可以通过扩展节点状态实现；
观点演化和信息扩散可以把邻居通道改为带方向或带符号的影响；
耦合振子、生态网络和电力系统则需要把离散分类读出替换为连续状态预测，
并采用与相互作用形式相符的聚合算子。

这种推广并非简单更换数据集。
不同动力学中的“结构证据”并不相同：SIR 中是感染事件与活跃感染者数量，振荡系统中可能是相位差和同步破缺，
生态系统中可能是扰动后的协同变化。
网络也可能有方向、权重、符号和时间变化。
因而，
本文提供的是一个方法框架和一组设计原则，
而不是无需修改即可覆盖所有复杂系统的通用成品模型。

\subsection{方法局限与实验边界}

模型的综合优势需要结合各指标和实验条件理解。
已知拓扑实验中，均值聚合模型的平均 $\Delta\gamma=0.00050$ 小于双通道的 0.00135；
ER 上三层求和 GraphConv 的自回归 MAE 为 0.0090，
也略低于双通道的 0.0106。
因此，
跨网络综合性能优势并不意味着每项指标均最优。
未知拓扑模型对 BA 的感染率反演值为 0.2889，
仍存在约 0.011 的低估。

九次优化中，ER 随机种子 7 和 BA 随机种子 21 分别生成 2955 条和 4586 条预测边，F1 为 0.2949 和 0.1193；
WS 随机种子 21 将 300 条真实边扩张为 390 条。
最低验证交叉熵在本组实验中排除了严重稠密解，
但三个初始化不足以估计稳定的失败概率，
也不能保证对新网络具有同样的筛选效果。

此外，未知拓扑阶段使用 100 轮训练和三个初始化选模，已知拓扑阶段使用 30 轮训练和单个初始化。
两阶段结果适合比较误差量级，不构成严格匹配训练预算的性能排名；
现有对照也未完全隔离参数量、采样与优化设置的影响。

本文的实验边界还体现在以下方面。
\begin{enumerate}
  \item {\yahei 网络与参数范围有限。}ER、WS、BA 各只有一个 $N=100$ 固定实例，且只使用 $\beta=0.3,\gamma=0.1$。
当前结果支持实例级方法验证，
不能替代跨网络、跨参数的统计结论。
  \item {\yahei 观测条件理想化。}节点身份已对齐，状态同步、完整且无标签噪声，网络静态、无向、无自环。
现实数据中的漏报、异步采样、潜变量和时变连接都会增加辨识难度。
  \item {\yahei 泛化检验仍较弱。}训练与验证按状态转移对随机划分，同一轨迹的相邻时刻可能进入不同集合；
因此验证交叉熵主要衡量同分布拟合与选模，
而非对全新轨迹的严格外推。
  \item {\yahei 联合优化仍不稳定。}三个训练 seed 只能展示初始化敏感性，不能估计可靠的失败概率。
生成器还需要维护平方级数量的候选关系，
对大规模网络的计算和搜索压力较高。
\end{enumerate}

这些限制并不否定已观察到的结果，但决定了结论应如何表达：本文证明了在受控 SIR 场景中，由状态时间序列联合恢复规律与结构是可行的，并找到了影响稳定性的关键机制；
它尚未证明任何满足局部交互的系统都能被同样准确地重构。

\subsection{未来研究方向}

未来研究将围绕“其他理论模型上的验证”与“真实复杂系统中的验证和拓展”两个方向展开，
逐步检验本文方法在不同动力学机制与观测条件下的适用性。

第一，在其他网络动力学理论模型上开展受控验证。
以复杂网络上的 Kuramoto 模型等为代表，
将研究对象从离散传播状态拓展到连续动力学，
根据相应的局部相互作用形式调整邻居表示、聚合算子和预测读出。
在已知真实规律与拓扑的条件下，
考察模型能否同时恢复局部交互、重构网络连接并预测群体演化。
通过多网络实例、不同参数条件和轨迹级划分，
进一步辨析双通道表示与交替结构学习的共性优势及其适用边界。

第二，在真实复杂系统中验证并拓展该框架。
面向金融市场、电网系统等场景，
探索如何从实际观测时间序列中学习组分间的有效相互作用，
并将所学关系用于动态预测和系统行为分析。
针对真实数据中的噪声、缺失、非平稳性和潜在外部驱动，
结合领域知识扩展有向、加权或时变的关系表示，
通过样本外预测、结构稳定性以及可获得的外部信息检验模型。
尤其需要区分预测性的依赖关系与真实物理连接或因果作用，
避免把较高预测精度直接解释为结构辨识成功。

这两个方向相互补充：理论模型提供可核验的方法检验基础，
真实系统则推动观测处理、关系表示与评价方式的进一步完善。
最终目标是形成能够连接时间序列观测、动力学预测与网络结构解释的复杂系统学习方法。
